\subsection{Motivación: combinar evidencia global y local}

Checkpoint~05 se define como una fase nueva porque combina explícitamente dos líneas
metodológicas que habían sido evaluadas por separado. Checkpoint~03 identificó representaciones
y modelos globales competitivos ---PLS, Ridge, CatBoost y ExtraTrees---. Checkpoint~04 mostró
que la geografía, la localidad y los grafos transductivos contienen señal adicional que no es
capturada completamente por esos modelos globales.

La hipótesis de Checkpoint~05 es, por tanto, una hipótesis de complementariedad de errores:
\[
\mathrm{Cov}(e^{Global},e^{Local})<\mathrm{Var}(e^{Local})
\]
en una magnitud suficiente para que una combinación validada reduzca RMSE respecto al
incumbente Local04D.

El checkpoint no supone que mayor complejidad implica mayor calidad. El incumbent permanece
\artifact{LocalRidge_C4_all_deterministic_Geo_k24_a30_p1}, con RMSE target-matched medio
0.487845, hasta que un procedimiento de selección split-excluded justifique reemplazarlo.

\subsection{Expertos base}

El conjunto de expertos globales hereda únicamente familias ya justificadas por Checkpoint~03:
\begin{align*}
G_1 &= \mathrm{PLS}(C0),\\
G_2 &= \mathrm{Ridge}(C3),\\
G_3 &= \mathrm{CatBoost}(C1),\\
G_4 &= \mathrm{ExtraTrees}(C1).
\end{align*}

PLS y Ridge representan modelos globales regularizados/latentes; CatBoost y ExtraTrees
representan respuestas no lineales de árboles. CatBoost se promedia sobre tres seeds después
de seleccionar hiperparámetros dentro del conjunto visible.

Los expertos transductivos se congelan desde 04D.1:
\[
L=\mathrm{LocalRidge}(C4,\mathrm{Geo},k=24,\alpha=30,p=1)
\]
y
\[
H=\mathrm{GraphDirect}(\mathrm{GeoAgro25},k=6,\lambda=8).
\]

Se reconstruyen también los ensambles históricos:
\[
E13=\frac{1}{2}(G_1+G_3),
\]
\[
E123=\frac{1}{3}(G_1+G_2+G_3),
\]
pero ahora dentro de las mismas pseudo-competiciones de Checkpoint~04. Esta diferencia es
crítica: no se mezclan OOF de 03C.3 con pseudo-predicciones generadas bajo particiones
distintas.

\subsection{Cross-fitting del nivel base}

Para una pseudo-competición outer con conjunto visible $\Train^\star$, el meta-modelo no puede
entrenarse sobre predicciones in-sample de los expertos. Se construye, por tanto, una matriz
OOF:
\[
Z_i=
\left[
\hat y_i^{G_1,-f(i)},
\hat y_i^{G_2,-f(i)},
\hat y_i^{G_3,-f(i)},
\hat y_i^{G_4,-f(i)},
\hat y_i^{L,-f(i)},
\hat y_i^{H,-f(i)}
\right].
\]

Cada columna es una predicción para $i$ obtenida sin utilizar $y_i$ en el fit correspondiente.
Los expertos globales realizan además selección de hiperparámetros dentro del training fold de
cada cross-fit. Para C3, cualquier expresión supervisada se vuelve a descubrir únicamente
dentro de dicho training fold.

La jerarquía es:
\[
\text{outer pseudo-test}
\supset
\text{meta cross-fit}
\supset
\text{inner tuning}.
\]
Aunque computacionalmente costosa, esta estructura evita que el meta-modelo aprenda del
optimismo in-sample de sus propios expertos.

\subsection{Descriptores de soporte para gating}

Para cada query se construye un vector de soporte $s_i$ relativo a las etiquetas actualmente
visibles:
\[
s_i=
(
\log(1+d_i^{geo}),
\log(1+d_i^{agro}),
q_i^{geo},
q_i^{agro},
r_i^{state},
r_i^{mun},
|\hat y_i^L-\hat y_i^H|
).
\]

$d_i^{geo}$ y $d_i^{agro}$ son distancias al vecino visible más cercano.
$q_i^{geo}$ y $q_i^{agro}$ expresan la posición de esas distancias respecto a la distribución
leave-one-out de vecinos entre las propias etiquetas visibles. $r_i^{state}$ y $r_i^{mun}$
son fracciones de etiquetas visibles del mismo estado/municipio. La discrepancia Local--Graph
es una señal de sensibilidad a inductive bias disponible en el momento de inferencia.

Estos descriptores no usan $y_i$ del query.

\subsection{Mixture-of-experts condicionado en soporte}

El meta-modelo principal es un gate softmax. Con $M$ expertos:
\begin{equation}
w_{im}
=
\frac{\exp(\eta_{im})}
{\sum_{\ell=1}^{M}\exp(\eta_{i\ell})},
\qquad
\eta_{im}=\beta_{m0}+\beta_m^\top s_i.
\label{eq:cp05-softmax-gate}
\end{equation}
La predicción es
\begin{equation}
\hat y_i
=
\sum_{m=1}^{M}w_{im}\hat y_{im},
\qquad
w_{im}\ge0,
\qquad
\sum_m w_{im}=1.
\label{eq:cp05-moe}
\end{equation}

El último experto se fija como logit de referencia para resolver la invariancia aditiva del
softmax. Los coeficientes del gate se regularizan:
\[
\mathcal L(\beta)
=
\frac{1}{n}\sum_i
(y_i-\hat y_i)^2
+
\lambda\|\beta\|_2^2.
\]
$\lambda$ se selecciona únicamente usando las filas meta-OOF.

\subsection{Controles de complejidad}

El mixture-of-experts se compara contra alternativas más simples:
\begin{itemize}
\item convex stacking con pesos globales no negativos que suman uno;
\item Ridge stacking;
\item Elastic Net stacking;
\item Huber stacking;
\item ExtraTrees shallow;
\item HistGradientBoosting shallow;
\item blends fijos Local+E13 y Local+E123.
\end{itemize}

Esto permite distinguir dos explicaciones. Si el mejor método es convex stacking, la ganancia
proviene de diversidad global de errores. Si MoE gana de forma reproducible, existe evidencia
adicional de que la calidad relativa de los expertos cambia con el soporte de la parcela.

\subsection{Protocolo de promoción}

Sea $m^\star$ el candidato con menor RMSE medio en la tabla completa de 16 pseudo-competiciones
target-matched. Ese mínimo no basta para reemplazar Local04D.

Para cada split $h$ se vuelve a seleccionar método usando los otros 15:
\[
\hat m_{-h}
=
\arg\min_{m\in\mathcal M}
\frac{1}{15}\sum_{s\neq h}\RMSE_{s,m}.
\]
La predicción LOSO concatenada debe mejorar simultáneamente:
\[
\overline{\RMSE}_{LOSO}
<
\overline{\RMSE}_{Local04D}
\]
y
\[
\RMSE_{\mathrm{pooled,LOSO}}
<
\RMSE_{\mathrm{pooled,Local04D}}.
\]

Además, $m^\star$ debe mejorar Local04D tanto en mean como pooled RMSE sobre la tabla de
desarrollo. Sólo entonces se promueve $m^\star$ al fit completo de las 138 etiquetas.

\subsection{Banco target-matched de confirmación fresca}

Existe una consideración adicional: los 16 splits target-matched originales ya participaron
en el refinamiento de Local04D durante 04D.1. Por ello, aunque el LOSO de Checkpoint~05
controla la selección de método dentro de ese banco, no constituye evidencia completamente
nueva respecto al incumbent.

Se predeclara entonces un segundo banco de 16 máscaras. Las máscaras se generan únicamente
con $X$ y metadata usando el mismo algoritmo de matching de 04B, pero con una seed no utilizada
durante el desarrollo. Si $\mathcal S_D$ denota el banco de desarrollo y $\mathcal S_C$ el
banco de confirmación, se impone al menos
\[
S_C\neq S_D
\qquad
\forall S_C\in\mathcal S_C,\ S_D\in\mathcal S_D,
\]
es decir, no se permiten máscaras exactamente duplicadas.

El challenger $m^\star$ se congela \emph{antes} de revelar cualquier $y$ de
$\mathcal S_C$. En confirmación sólo se compara
\[
\{m^\star,\mathrm{Local04D}\}.
\]
La promoción exige simultáneamente
\[
\overline{\RMSE}_{C,m^\star}
<
\overline{\RMSE}_{C,Local}
\]
y
\[
\RMSE_{\mathrm{pooled},C,m^\star}
<
\RMSE_{\mathrm{pooled},C,Local}.
\]
Si el gate de desarrollo falla, el banco de confirmación ni siquiera se utiliza para escoger
otro método. Esto preserva su función como prueba confirmatoria y evita convertirlo en un
segundo search space.

\subsection{Exportación}

El runner produce el archivo de predicciones final y dos bundles serializados:
\[
\artifact{models/final/checkpoint05_app_bundle.joblib},
\qquad
\artifact{models/final/checkpoint05_model.joblib}.
\]
El primero es deliberadamente ligero: contiene la predicción congelada de las 59 parcelas,
candidatos, diagnósticos, manifest y esquema, sin objetos entrenados de CatBoost o
scikit-learn. Es el artefacto preferido para la futura app cuando sólo se necesita reproducir
el conjunto fijo de competencia.

El segundo conserva además los expertos globales ajustados y los meta-modelos para
reproducibilidad del procedimiento. En esta fase ambos tienen scope de inferencia
explícitamente restringido a los 59 objetivos fijos. Una API para parcelas arbitrariamente
nuevas se implementará separadamente para no confundir ambos contratos.

\subsection{Resultado ejecutado y decisión}

El run final completó el 23 de septiembre de 2026. El mejor challenger de la tabla completa
de desarrollo fue:
\[
\mathrm{LocalE13}_{0.10}
=
0.90\,\mathrm{Local04D}
+
0.10\,E13.
\]
Sus métricas fueron
\[
\overline{\RMSE}=0.487243,
\qquad
\RMSE_{\mathrm{pooled}}=0.495261,
\]
frente a Local04D:
\[
\overline{\RMSE}=0.487845,
\qquad
\RMSE_{\mathrm{pooled}}=0.495741.
\]
La ventaja nominal fue menor a una milésima.

El gate LOSO no confirmó la mejora:
\[
\overline{\RMSE}_{LOSO}=0.487934,
\qquad
\RMSE_{\mathrm{pooled,LOSO}}=0.495933,
\]
ligeramente peor que Local04D en ambas métricas requeridas. Por contrato, el banco fresco de
confirmación no se puntuó para buscar una alternativa distinta.

Una ejecución previa se detuvo en 29/34 outer splits porque un fold nested de 18 observaciones
hacía inviable PLS con 20 componentes. Se corrigió únicamente la factibilidad:
\[
K_{\max}
=
\min\left(p,\min_f n_{\mathrm{train},f}\right),
\]
eliminando candidatos con $n_{\mathrm{components}}>K_{\max}$ dentro de cada fit. La regla no
usa $y$ para seleccionar complejidad y permitió reutilizar de forma segura los splits
completados.

No hubo promoción. La fuente canónica posterior a Checkpoint~05 es
\begin{center}
\artifact{reports/checkpoint_05/final_predictions.csv},
\end{center}
que conserva exactamente Local04D para las 59 parcelas.

Checkpoint~03D se ejecutó posteriormente como cierre de capacidad global. Sus scores
state/grouped no retroactivan una promoción en 05 porque pertenecen a protocolos distintos.
Cualquier nueva mezcla con esos globales requiere un 05B explícito.
